%!TeX root=../main.tex

\chapter{Proposta}
\label{cap:proposta}


\section{Câmera}
\label{sec:camera}

A visualização do ambiente de jogo adota uma perspectiva tridimensional com
projeção ortográfica, configurada sob uma inclinação espacial de estilo
isométrico, focando na legibilidade espacial dos cenários.

Para auxiliar nesta legibilidade, diferente de abordagens tradicionais com ângulos estáticos de visualização,
o sistema de câmera do projeto implementa um grau de liberdade de rotação
horizontal contínua em $360^\circ$ ao redor do personagem principal. Essa rotação
permite ao jogador girar o ponto de vista livremente para examinar cantos
ocultos do mapa, antecipar ameaças e navegar por rotas complexas sem
alterar a inclinação vertical fixa da projeção. Esta técnica é comumente
encontrada em CRPGs (\textit{Computer Role-Playing Games}) isométricos modernos, como \textit{Divinity: Original Sin 2} (\cite{original-sin}).

Para assegurar o ritmo acelerado (\textit{fast-paced}) da jogabilidade, o
acompanhamento da câmera em relação ao movimento do jogador é projetado
para ocorrer de forma suave (\textit{smooth follow}). Esse amortecimento visual
evita saltos bruscos resultantes de mudanças de direção repentinas do
personagem, mantendo a cena estável.

Por fim, o comportamento da câmera foi especificado para lidar com os
desafios de oclusão inerentes à perspectiva isométrica. Ao posicionar o
jogador atrás de estruturas altas ou elementos volumosos do cenário, a
câmera e os sistemas visuais integrados devem garantir a manutenção do
\textit{feedback} de localização do personagem, preservando a orientação espacial e
a fluidez dos comandos em qualquer direção do mapa.

\section{Movimentação}
\label{sec:movimentacao}

O sistema de locomoção do personagem principal foi estruturado para
atender aos requisitos de alta velocidade e resposta imediata inerentes do
gênero \textit{fast-paced}. O controle de deslocamento ocorre de maneira
omnidirecional no plano horizontal do cenário, permitindo ao jogador mover-se
fluidamente em qualquer direção independentemente da rotação atual da
câmera.

Para garantir o dinamismo exigido pelo projeto, o conjunto de ações
fundamentais de movimento é composto em:

\begin{description}
	\item[Deslocamento:] O personagem responde de forma
	contínua aos comandos direcionais.
	\item[Corrida:] O jogador dispõe de um estado de corrida
	que aumenta a sua velocidade de locomoção pelo mapa.
	\item[Pulo:] Mecânica de salto para permitir a navegação vertical
	sobre as plataformas sem quebrar o fluxo contínuo de movimento.
\end{description}

Do ponto de vista de \textit{game design}, as transições entre estados de movimento
(como a entrada e saída da corrida ou a execução de saltos) foram desenvolvidas
com foco na conservação de momento. Em vez de travamentos ou freadas abruptas,
o sistema de movimento preserva a energia cinética acumulada pelo jogador. Isto
permite, por exemplo, que a velocidade obtida durante uma corrida ou um deslize
seja transferida para o pulo. Essa continuidade física assegura que o jogador
mantenha total controle e fluidez ao se locomover pelo ambiente.


\section{Progressão e Mecânicas de Jogo}
\label{sec:progressao_e_mecanicas}


\section{Efeitos Visuais}
\label{sec:efeitos_visuais}


\section{Pós-processamento}
\label{sec:pos_processamento}


\section{Interação do Usuário}
\label{sec:interacao_usuario}


\section{Som}
\label{sec:som}


\section{Configurações}
\label{sec:configuracoes}